\section{Numerical methods}
\label{sec:numerics}

\subsection{Solving the closure}

\Cref{eq:closure-q} could be solved by radicals at $q=3/2$ or by Newton's
method in general. Neither is used, for reasons that are structural rather than
matters of taste.

Newton's method applied to $F(m)-\Mzero$ has a critical defect here: by
\cref{lem:concave}, $F'(m^*)=0$, so the Newton step is unbounded near the fold,
which is precisely the region of interest. Worse, from a starting point on the
wrong side it converges to $m_+$, the unstable branch, without any indication
that it has done so.

The engine instead exploits the shape established in \cref{lem:concave}:

\begin{enumerate}[nosep]
  \item Compute $m^*$ and $\Mmax$ in closed form from
        \cref{eq:mstar,eq:Mmax}. Feasibility is decided \emph{analytically} by
        comparing $\Mzero$ with $\Mmax$, never by whether an iteration
        converged.
  \item If feasible, bracket $m_-$ on $(0,m^*]$, where $F$ is a strictly
        increasing bijection, and bisect. Bracket $m_+$ on $[m^*,\infty)$ by
        geometric expansion and bisect.
\end{enumerate}

Bisection converges linearly with rate $1/2$ and requires
$\lceil\log_2((b-a)/\epsilon)\rceil$ evaluations, about $40$ for double
precision. It is slower than Newton and it cannot fail. Given that a single
evaluation of $F$ costs microseconds, and given that the failure modes of
Newton here are silent rather than loud, this is the correct trade.

\begin{remark}[Feasibility is analytic, not numerical]
This is the single most important implementation decision in the sizing layer,
and it follows from the critical slowing-down remark after
\cref{thm:iteration}. A solver that reports non-convergence after $K$ iterations
cannot distinguish a near-feasible design from an infeasible one, because
$\mathrm{rate} = 1-F'(m_-)\to1$ as $\Mzero\uparrow\Mmax$. Deciding feasibility
by $\Mzero\lessgtr\Mmax$ removes the ambiguity entirely.
\end{remark}

For the component closure \cref{eq:fpclosure} no closed form exists, and the
shape of $\mathcal{F}$ is observed rather than proved (\cref{sec:components}),
so the search cannot borrow the analytic feasibility test. The engine samples
$\mathcal{F}$ on a geometric grid of ratio $1.35$ from \SI{e-3}{\kilo\gram}
upward, stopping at the first non-negative sample. Every sampled local maximum
is then refined by bounded minimisation of $-\mathcal{F}$ in $\ln m$, because a
positive stretch of $\mathcal{F}$ narrower than one grid interval can sit
between two negative samples. The first bracket with a sign change, in
increasing mass, is closed by Brent's method to a relative tolerance of
$10^{-12}$, which selects the light branch. The search reports one of four
outcomes: a root, a near-tangent root resolved only to residual tolerance, a
non-finite residual, or an exhausted search. The last is reported as a failure
of a bounded search and is never described as infeasibility. The case that
motivated the refinement is kept as a regression test: the design of
\cref{sec:results} at $t_f=\SI{4280.0502}{\second}$ has a positive island of
width \SI{5}{\percent} in mass, with roots at \SI{4.5678}{\kilo\gram} and
\SI{4.6173}{\kilo\gram}, which the grid alone straddled and reported as no
mass closing.

\subsection{The battery loop with the simulator inside it}

\Cref{eq:fixedpoint} is solved by iteration, but naive iteration converges only
at the linear rate $\mathcal{H}'(\mu)$ of \cref{prop:contraction}, which near
the fold approaches unity. Since each evaluation of $\mathcal{H}$ costs a full
trajectory simulation, this is expensive.

The engine applies a secant step to the residual
$\mathrm{gap}(\mu) = \mathcal{H}(\mu)-\mu$:
\begin{equation}
  \mu_{k+1} = \mu_k - \mathrm{gap}(\mu_k)\,
    \frac{\mu_k-\mu_{k-1}}{\mathrm{gap}(\mu_k)-\mathrm{gap}(\mu_{k-1})} ,
  \label{eq:secant}
\end{equation}
which uses only quantities already computed and therefore costs no additional
simulations.

\begin{proposition}[Convergence order]
\label{prop:secant}
If $\mathrm{gap}\in C^2$ near a simple root $\mu^\star$ with
$\mathrm{gap}'(\mu^\star)\ne0$, the secant iteration \cref{eq:secant} converges
locally with order $\varphi=(1+\sqrt5)/2\approx1.618$, against order $1$ for
the fixed-point iteration.
\end{proposition}

\begin{proof}
Classical; see Ortega and Rheinboldt~\cite[\S8.2]{ortega1970}. The error
satisfies $|\epsilon_{k+1}|\approx C|\epsilon_k||\epsilon_{k-1}|$, whose
characteristic exponent is the golden ratio.
\end{proof}

The step is guarded: it is accepted only when it moves in the same direction as
the gap and by no more than six times its magnitude, otherwise the plain
fixed-point step is taken. Because $\mathcal{H}$ takes the larger of an
energy and a power requirement (\cref{eq:outermap}) it has a kink where the
two exchange, and \cref{prop:secant} does not apply across it; the guard is
what keeps the iteration sane there. When the gap is within tolerance the
battery is rounded up to the last requirement and the design is flown once
more, and the loop reports convergence only if that flight requires no more
battery than it carries. On every mission of \cref{sec:results} this took
three updates and two verification flights, five simulations in all, against
eight or more for plain iteration.

\subsection{Trajectory integration}

\Cref{eq:ode} is integrated with classical fourth-order Runge--Kutta at fixed
step. Fixed step rather than adaptive, because the cost must be predictable for
an interactive tool and because the controller is evaluated once per step,
making the scheme effectively a sampled-data system whose sample rate should
not wander.

Local truncation error is $O(h^5)$ and global error $O(h^4)$. The step is
chosen from the fastest closed-loop mode. With $\tau_m=\SI{45}{\milli\second}$,
$K_R\approx\SI{41}{\per\second\squared}$ and hence
$\omega_n=\sqrt{K_R}\approx\SI{6.4}{\radian\per\second}$, the fastest pole is
the motor at $1/\tau_m\approx\SI{22}{\radian\per\second}$; at
$h=\SI{4}{\milli\second}$ the product $h\lambda\approx0.09$ sits deep inside the
RK4 stability region, whose real-axis extent is $|h\lambda|\le2.79$.

The quaternion is renormalised after each step, $q\leftarrow q/\|q\|$. This
does not cost the order. The exact flow preserves $\|q\|=1$, so the RK4 step
returns $\tilde q$ with $\|\tilde q\| = 1+O(h^5)$; the map
$\tilde q\mapsto\tilde q/\|\tilde q\|$ is smooth near $\mathbb{S}^3$ with
$\tilde q/\|\tilde q\| = \tilde q + O(\|\tilde q\|-1)$, so the projection
perturbs each step by $O(h^5)$, the same order as the local truncation error,
and the global order $O(h^4)$ is retained.

Two different things are being integrated and the distinction matters. The
controller is evaluated once per step and its output held, so RK4 integrates
the plant under a piecewise-constant input at fourth order. The closed loop is
therefore a sampled-data system at \SI{250}{\hertz}, which is what a digital
autopilot is; the $O(h)$ difference between it and a continuous-feedback
idealisation is a property of the system being simulated, not an integration
error. A convergence check with the controller rate held fixed and the
integration substeps refined was not run and is listed in
\cref{sec:limitations}.

\begin{remark}[Windowed mission energy]
Integrating a full \SI{1800}{\second} mission at $h=\SI{4}{\milli\second}$ would
require $4.5\times10^5$ steps per evaluation, which is incompatible with the
outer loop. The engine instead integrates a representative window of
\SIrange{12}{20}{\second}, measures the mean electrical power after a settling
interval, and extrapolates:
$E_{\mathrm{mission}} = \bar P\,t_f$. This is exact for a periodic mission whose
window spans a whole number of periods and is otherwise an approximation whose
error is the difference between the window mean and the mission mean. It is
recorded as a limitation.
\end{remark}

\begin{remark}[Initial conditions]
The initial state is taken to be the reference position \emph{and the reference
velocity}. Starting at rest beside a person already walking is not a gentler
initial condition but a worse one: on a curved path the display sits in the
person's path, so they close on it while it accelerates. This artefact
contaminated every mission statistic until it was found; see
\cref{sec:verification}.
\end{remark}

\subsection{Linearisation}

$\mat A$ and $\mat B$ are formed by central differences,
\begin{equation}
  \mat A_{:,j} \approx \frac{f(x_0+h\vect e_j)-f(x_0-h\vect e_j)}{2h},
  \qquad h = \varepsilon\max\{1,|x_{0,j}|\},
  \label{eq:fd}
\end{equation}
with $\varepsilon=10^{-6}$. Central differences have truncation error $O(h^2)$
against $O(h)$ for one-sided, and the total error is minimised near
$h\sim\epsilon_{\mathrm{mach}}^{1/3}\approx6\times10^{-6}$, so the choice is
near-optimal. Controllability is assessed by the reachable-subspace iteration
of \cref{sec:control} with a singular-value tolerance of $10^{-9}$ on the
norm-scaled matrices, and requires full rank.

\subsection{The outer optimisation}

\Cref{eq:outerprog} is solved by differential evolution~\cite{storn1997}, a
population-based stochastic method, with constraints handled by an exterior
quadratic penalty
\begin{equation}
  \tilde{\mathcal{J}}(p) = \mathcal{J}(p)
    + \varpi\sum_j \left(\frac{\max\{0,g_j(p)\}}{s_j}\right)^{2} ,
  \label{eq:penalty}
\end{equation}
with per-constraint scales $s_j$ chosen so that a violation of engineering
significance contributes order unity.

Three properties of the problem dictate this choice. The objective is
non-convex and its landscape contains large regions where the closure has no
solution at all, so gradient methods have nowhere to start. Several design
variables are integer (rotor count, blade count), so smoothness fails. And an
evaluation of the algebraic objective costs milliseconds, so a population of a
few hundred is affordable.

The penalty is exterior rather than a barrier precisely because the infeasible
region must remain traversable: a barrier method cannot cross from one feasible
island to another through infeasible territory, and here the islands are
separated by the fold.

\subsection{Numerical hygiene}

Three details are worth recording because each caused a visible defect.

\emph{Non-finite values in transport.} \texttt{NaN} and $\pm\infty$ are not
representable in JSON. Infeasible cells in a sweep, and quantities such as
$\Mmax$ when no fold exists, are genuinely undefined rather than zero. The API
maps every non-finite float to \texttt{null} and the front end renders it as a
dash, so an undefined quantity is never silently displayed as $0$.

\emph{Solidity and Reynolds guards.} Without \cref{as:sigma} and
\cref{eq:recorrection} the optimiser drives solidity to $0.55$ and tip speed
toward zero, exploiting the fact that \cref{prop:profile} has no notion of
blade interference and \cref{prop:fmblade} none of Reynolds number. Both
constraints exist to keep the optimiser inside the domain where the model is a
model of something.

\emph{Boundary comparisons.} Constraints active exactly at their bound were
being reported as violations because of floating-point comparison. All
constraint checks carry a relative tolerance of $2\times10^{-3}$, chosen so
that a design sitting on a bound reads as satisfying it.
